%%%%%%%%%%%%%%%%%%%%%%%%%%%%%%%%%%%%%%%%%%%%%%%%%%%%%%%%%%%%%%%%%%%%%%%%%%%%%%%
%  SÍLABO DE ECONOMÍA MATEMÁTICA — Plan de estudio integrado
%
%  Textos base:
%    CG  — Chávez y Gallardo, Álgebra lineal y optimización para el análisis
%          económico (Fondo Editorial PUCP, 2026) — ejemplar físico; se cita
%          por el índice de capítulos y secciones
%    CT  — Carter, Foundations of Mathematical Economics (MIT Press, 2001)
%    DF  — De la Fuente, Mathematical Methods and Models for Economists
%          (Cambridge University Press, 2000)
%  Optimización dinámica:
%    KS  — Kamien y Schwartz, Dynamic Optimization (Elsevier, 2.ª ed., 1991)
%    LR  — Lomelí y Rumbos, Métodos dinámicos en economía (ITAM)
%    MS  — Meinsma y van der Schaft, A Course on Optimal Control (Springer, 2023)
%    DP  — Sargent y Stachurski, Dynamic Programming, vol. I (2025)
%    ST  — Stachurski, Economic Dynamics: Theory and Computation (MIT, 2.ª ed.)
%    BL  — Bonifaz y Lama, Optimización dinámica y teoría económica (UP)
%    LU  — Luenberger, Optimization by Vector Space Methods (Wiley, 1969)
%    EL  — Elaydi, An Introduction to Difference Equations (Springer, 3.ª ed.)
%  Programas PUCP integrados: MAT291 (Matemáticas para economistas),
%    1ECO75 (Economía Matemática Intermedia), MAT693 (Concavidad y
%    Optimización), MAT659 (Temas Avanzados en Economía Matemática) y
%    MAT835 (Control Óptimo Estocástico); de referencia, MAT616 y MAT656.
%
%  Compilar con:  pdflatex EconomiaMatematica.tex   (dos veces, por el índice)
%%%%%%%%%%%%%%%%%%%%%%%%%%%%%%%%%%%%%%%%%%%%%%%%%%%%%%%%%%%%%%%%%%%%%%%%%%%%%%%

\documentclass[11pt,a4paper]{article}

\usepackage[T1]{fontenc}
\usepackage[utf8]{inputenc}
\usepackage[spanish,es-noquoting,es-tabla]{babel}
\usepackage{lmodern}
\usepackage[scaled=0.92]{helvet}        % títulos en sans serif, al estilo McGraw-Hill
\usepackage[margin=2.4cm,headheight=15pt]{geometry}
\usepackage[table]{xcolor}
\usepackage{enumitem}
\usepackage{titlesec}
\usepackage{titletoc}
\usepackage{tabularx}
\usepackage{booktabs}
\usepackage{longtable}
\usepackage{amsmath}
\usepackage{amssymb}
\usepackage{fancyhdr}
\usepackage{tikz}
\usepackage[skins]{tcolorbox}
\usepackage{needspace}

%------------------------------------------------------------------ colores ---
% Gama de violetas
\definecolor{tonoOscuro}{RGB}{46,4,82}              % oscuro
\definecolor{tonoProfundo}{RGB}{68,6,121}           % profundo
\definecolor{tonoPrincipal}{RGB}{89,8,160}          % principal
\definecolor{tonoA}{RGB}{138,25,179}                % variante cálida
\definecolor{tonoB}{RGB}{75,15,179}                 % variante fría
\definecolor{tonoC}{RGB}{56,21,193}                 % variante viva
\definecolor{tonoClaro}{RGB}{162,95,221}            % claro
\definecolor{tonoAcero}{RGB}{99,63,131}             % acero
\definecolor{tonoPizarra}{RGB}{111,86,134}          % pizarra
\definecolor{tonoGrafito}{RGB}{57,39,73}            % grafito (texto secundario)
\definecolor{tonoHielo}{RGB}{228,207,247}           % hielo
\definecolor{tonoPalido}{RGB}{243,233,251}          % pálido
\definecolor{tonoBruma}{RGB}{250,246,254}           % bruma

\colorlet{azulosc}{tonoOscuro}
\colorlet{grisclaro}{tonoPalido}
% Un tono distinto de la gama para cada libro
\definecolor{colorCG}{RGB}{36,15,138}
\definecolor{colorCT}{RGB}{60,20,184}
\definecolor{colorDF}{RGB}{63,18,161}
\definecolor{colorKS}{RGB}{63,15,138}
\definecolor{colorLR}{RGB}{95,20,184}
\definecolor{colorMS}{RGB}{94,18,161}
\definecolor{colorDP}{RGB}{90,15,138}
\definecolor{colorST}{RGB}{131,20,184}
\definecolor{colorBL}{RGB}{125,18,161}
\definecolor{colorLU}{RGB}{116,15,138}
\definecolor{colorEL}{RGB}{167,20,184}
\colorlet{colorPUCP}{tonoPizarra}

\usepackage[colorlinks=true,linkcolor=tonoOscuro,urlcolor=tonoPrincipal]{hyperref}

%------------------------------------------------------------------ estilos ---
\titleformat{\section}[block]
  {\sffamily\LARGE\bfseries\color{tonoOscuro}\raggedright}{\thesection.}{0.6em}
  {\textcolor{tonoPrincipal}{\rule[-0.1ex]{0.55em}{0.85em}}\hspace{0.45em}}
  [\vspace{-0.4ex}{\color{tonoPrincipal}\titlerule[1.6pt]}\vspace{1.2pt}{\color{tonoClaro}\titlerule[0.6pt]}]
\titlespacing*{\section}{0pt}{2.6ex plus 1ex minus .2ex}{1.6ex plus .2ex}

\titleformat{\subsection}
  {\sffamily\large\bfseries\color{tonoPrincipal}}{\thesubsection}{0.6em}
  {\textcolor{tonoClaro}{\rule[0.1ex]{0.45em}{0.45em}}\hspace{0.45em}}
\titlespacing*{\subsection}{0pt}{2.2ex plus .8ex minus .2ex}{1ex plus .2ex}

% Índice
\titlecontents{section}[0pt]
  {\addvspace{3pt}\sffamily\bfseries\color{tonoOscuro}}
  {}{\textcolor{tonoB}{\rule[0.15ex]{0.4em}{0.4em}}\hspace{0.6em}}
  {{\color{tonoClaro}\titlerule*[0.7em]{.}}\textcolor{tonoPrincipal}{\contentspage}}

% Cabecera y pie
\pagestyle{fancy}
\fancyhf{}
\lhead{\small\sffamily\bfseries\color{tonoOscuro}Sílabo de Economía Matemática}
\rhead{\footnotesize\sffamily\color{tonoPrincipal}Chávez--Gallardo $\cdot$ Carter $\cdot$ De la Fuente $\cdot$ Kamien--Schwartz}
\cfoot{{\setlength{\fboxsep}{3pt}\colorbox{tonoPrincipal}{\makebox[1.8em]{\small\sffamily\bfseries\color{white}\thepage}}}}
\renewcommand{\headrule}{%
  {\color{tonoPrincipal}\hrule height 1.1pt width\headwidth}\vskip 1.2pt
  {\color{tonoClaro}\hrule height 0.4pt width\headwidth}\vskip -2.7pt}

% Listas
\setlist[itemize,1]{label=\textcolor{tonoPrincipal}{\scriptsize$\blacksquare$}}
\setlist[enumerate,1]{label=\textcolor{tonoPrincipal}{\sffamily\bfseries\arabic*.}}

% Tablas: cabecera oscura y filas alternas en tono pálido
\arrayrulecolor{tonoPrincipal}
\setlength{\aboverulesep}{0pt}
\setlength{\belowrulesep}{0pt}
\renewcommand{\arraystretch}{1.3}
\newcommand{\cab}[1]{\textcolor{white}{\sffamily\bfseries #1}}
\newcommand{\filasalternas}{\rowcolors{2}{white}{tonoPalido}}

% Encabezados de párrafo en color
\newcommand{\destacado}[1]{{\sffamily\bfseries\color{tonoPrincipal}#1}}

% Franja decorativa con la gama
\newcommand{\franjas}[1][3pt]{%
  \noindent\begin{tikzpicture}
    \foreach \c [count=\i from 0] in {tonoOscuro,tonoProfundo,tonoPrincipal,tonoA,tonoB,tonoC,tonoClaro,tonoHielo}
      \fill[\c] (\i*\linewidth/8,0) rectangle ++(\linewidth/8,#1);
  \end{tikzpicture}\par}

%---------------------------------------------------------------- comandos ----
% Etiqueta de libro: recuadro de color con la sigla
\newcommand{\etiqueta}[2]{%
  {\setlength{\fboxsep}{1.4pt}\colorbox{#1}{\sffamily\bfseries\color{white}\strut#2}}}
% Referencias abreviadas a cada libro
\newcommand{\CG}[1]{\etiqueta{colorCG}{CG}\nobreak\hspace{0.25em}\textcolor{colorCG}{#1}}
\newcommand{\CT}[1]{\etiqueta{colorCT}{CT}\nobreak\hspace{0.25em}\textcolor{colorCT}{#1}}
\newcommand{\DF}[1]{\etiqueta{colorDF}{DF}\nobreak\hspace{0.25em}\textcolor{colorDF}{#1}}
\newcommand{\KS}[1]{\etiqueta{colorKS}{KS}\nobreak\hspace{0.25em}\textcolor{colorKS}{#1}}
\newcommand{\LR}[1]{\etiqueta{colorLR}{LR}\nobreak\hspace{0.25em}\textcolor{colorLR}{#1}}
\newcommand{\MS}[1]{\etiqueta{colorMS}{MS}\nobreak\hspace{0.25em}\textcolor{colorMS}{#1}}
\newcommand{\DP}[1]{\etiqueta{colorDP}{DP}\nobreak\hspace{0.25em}\textcolor{colorDP}{#1}}
\newcommand{\ST}[1]{\etiqueta{colorST}{ST}\nobreak\hspace{0.25em}\textcolor{colorST}{#1}}
\newcommand{\BL}[1]{\etiqueta{colorBL}{BL}\nobreak\hspace{0.25em}\textcolor{colorBL}{#1}}
\newcommand{\LU}[1]{\etiqueta{colorLU}{LU}\nobreak\hspace{0.25em}\textcolor{colorLU}{#1}}
\newcommand{\EL}[1]{\etiqueta{colorEL}{EL}\nobreak\hspace{0.25em}\textcolor{colorEL}{#1}}

% Un tema: número en recuadro de color + título
\newcommand{\tema}[1]{\temaaux#1\relax}
\def\temaaux#1 #2\relax{%
  \par\addvspace{1.1em}\noindent
  {\sffamily\bfseries\large\raggedright
   \hangindent=2.9em\hangafter=1
   \makebox[2.9em][l]{{\setlength{\fboxsep}{2.5pt}\colorbox{tonoPrincipal}{\color{white}#1}}}%
   {\color{tonoOscuro}#2}\par}%
  \nobreak\vspace{0.25em}\noindent{\color{tonoClaro}\rule{\linewidth}{0.6pt}}\par\nobreak\vspace{0.35em}}

% Línea de lecturas con marcador y etiqueta de color
\newcommand{\linealectura}[3]{%
  \par\nobreak\noindent\hspace{1.2em}%
  \begin{minipage}{0.92\textwidth}\small
    \textcolor{#1}{\rule[0.1ex]{0.45em}{0.45em}}\hspace{0.45em}%
    {\sffamily\bfseries\color{#1}#2}\enspace#3%
  \end{minipage}\par\vspace{0.3em}}
\newcommand{\lecturas}[1]{\linealectura{tonoPrincipal}{Lecturas:}{#1}}
\newcommand{\apoyo}[1]{\linealectura{tonoA}{Profundización:}{#1}}
\newcommand{\juegos}[1]{\linealectura{tonoB}{Teoría de juegos:}{#1}}
\newcommand{\supl}[1]{\linealectura{tonoC}{Suplemento:}{#1}}
\newtcolorbox{cajanota}{enhanced, sharp corners, boxrule=0pt, frame hidden,
  colback=tonoPalido, borderline west={2.5pt}{0pt}{tonoClaro},
  left=8pt, right=6pt, top=4pt, bottom=4pt, before skip=0pt, after skip=0pt,
  fontupper=\small\color{tonoGrafito}}
\newcommand{\nota}[1]{%
  \par\nobreak\noindent\hspace{1.2em}%
  \begin{minipage}{0.92\textwidth}\begin{cajanota}#1\end{cajanota}\end{minipage}\par\vspace{0.3em}}
\newcommand{\consultalib}[1]{\linealectura{tonoPizarra}{Consulta:}{{\color{tonoGrafito}#1}}}
% Artículos y textos tomados de los programas PUCP
\newcommand{\articulos}[1]{\linealectura{colorPUCP}{Artículos (PUCP):}{#1}}
\newcommand{\consulta}[1]{\linealectura{colorPUCP}{Consulta (PUCP):}{{\color{tonoGrafito}#1}}}

% Encabezado de bloque: banda oscura, franja de la gama y recuadro de presentación
\newif\ifrecienparte
\newcommand{\bloque}[2]{%
  \ifrecienparte\global\recienpartefalse\else\par\needspace{0.4\textheight}\fi
  \phantomsection\addcontentsline{toc}{section}{#1}%
  \begin{tcolorbox}[enhanced, sharp corners, boxrule=0pt, frame hidden,
      colback=tonoOscuro, borderline west={7pt}{0pt}{tonoClaro},
      left=14pt, right=10pt, top=10pt, bottom=10pt, before skip=0pt, after skip=0pt]
    {\raggedright\sffamily\Large\bfseries\color{white}#1\par}
  \end{tcolorbox}%
  \nointerlineskip\franjas[4pt]
  \vspace{0.8em}
  \begin{tcolorbox}[enhanced, sharp corners, boxrule=0pt, frame hidden,
      colback=tonoPalido, borderline west={3pt}{0pt}{tonoPrincipal},
      borderline south={0.6pt}{0pt}{tonoClaro},
      left=10pt, right=8pt, top=6pt, bottom=6pt, before skip=0pt, after skip=0pt,
      fontupper=\small\color{tonoGrafito}]
    #2
  \end{tcolorbox}%
  \par\medskip}

% Encabezado de parte: banda a todo el ancho en el tono principal
\newcommand{\parte}[1]{%
  \clearpage
  \phantomsection\addcontentsline{toc}{section}{\textbf{#1}}%
  \begin{tcolorbox}[enhanced, sharp corners, boxrule=0pt, frame hidden,
      colback=tonoPrincipal, halign=center, top=16pt, bottom=16pt,
      borderline north={3pt}{0pt}{tonoClaro}, borderline south={3pt}{0pt}{tonoClaro},
      before skip=0pt, after skip=0pt]
    {\sffamily\LARGE\bfseries\color{white}#1}
  \end{tcolorbox}\par\bigskip
  \global\recienpartetrue}

%================================================================== documento ==
\begin{document}

\begin{titlepage}
\begin{tikzpicture}[remember picture,overlay]
  % banda superior
  \fill[tonoOscuro] (current page.north west) rectangle ([yshift=-5.4cm]current page.north east);
  \fill[tonoClaro] (current page.north west) rectangle ([xshift=1.1cm,yshift=-5.4cm]current page.north west);
  % franja de la gama bajo la banda
  \foreach \c [count=\i from 0] in {tonoOscuro,tonoProfundo,tonoPrincipal,tonoA,tonoB,tonoC,tonoClaro,tonoHielo}
    \fill[\c] ([xshift=\i*\paperwidth/8,yshift=-5.4cm]current page.north west)
      rectangle ++(\paperwidth/8,-0.35cm);
  % franja inferior
  \foreach \c [count=\i from 0] in {tonoHielo,tonoClaro,tonoC,tonoB,tonoA,tonoPrincipal,tonoProfundo,tonoOscuro}
    \fill[\c] ([xshift=\i*\paperwidth/8]current page.south west)
      rectangle ++(\paperwidth/8,0.6cm);
\end{tikzpicture}%
\centering
\vspace*{0.5cm}
{\sffamily\Huge\bfseries\color{white} Sílabo de Economía Matemática\par}
\vspace{0.4cm}
{\sffamily\Large\color{tonoHielo} Plan de estudio integrado\par}
\vspace{1.9cm}
{\sffamily\large\bfseries\color{tonoPrincipal}\textcolor{tonoClaro}{\rule[0.5ex]{1.5cm}{0.6pt}}\quad Textos base\quad\textcolor{tonoClaro}{\rule[0.5ex]{1.5cm}{0.6pt}}\par}
\vspace{0.3cm}
{\large\etiqueta{colorCG}{CG} \quad {\color{tonoOscuro}Jorge Chávez y Marcelo Gallardo}\par}
{\normalsize \emph{\color{colorCG}Álgebra lineal y optimización para el análisis económico}\par}
{\small\color{tonoPizarra} Fondo Editorial PUCP, 2026 --- 10 capítulos (ejemplar físico)\par}
\vspace{0.2cm}
{\large\etiqueta{colorCT}{CT} \quad {\color{tonoOscuro}Michael Carter}\par}
{\normalsize \emph{\color{colorCT}Foundations of Mathematical Economics}\par}
{\small\color{tonoPizarra} The MIT Press, 2001 --- 6 capítulos\par}
\vspace{0.2cm}
{\large\etiqueta{colorDF}{DF} \quad {\color{tonoOscuro}Ángel de la Fuente}\par}
{\normalsize \emph{\color{colorDF}Mathematical Methods and Models for Economists}\par}
{\small\color{tonoPizarra} Cambridge University Press, 2000 --- 13 capítulos\par}
\vspace{0.6cm}
{\sffamily\large\bfseries\color{tonoPrincipal}\textcolor{tonoClaro}{\rule[0.5ex]{1.5cm}{0.6pt}}\quad Optimización dinámica\quad\textcolor{tonoClaro}{\rule[0.5ex]{1.5cm}{0.6pt}}\par}
\vspace{0.3cm}
{\large\etiqueta{colorKS}{KS} \quad {\color{tonoOscuro}Morton I. Kamien y Nancy L. Schwartz}\par}
{\normalsize \emph{\color{colorKS}Dynamic Optimization: The Calculus of Variations and Optimal Control in Economics and Management}\par}
{\small\color{tonoPizarra} Elsevier, 2.ª edición, 1991 --- dos partes\par}
\vspace{0.2cm}
{\large\etiqueta{colorLR}{LR} \quad {\color{tonoOscuro}Héctor Lomelí y Beatriz Rumbos}\par}
{\normalsize \emph{\color{colorLR}Métodos dinámicos en economía}\par}
{\small\color{tonoPizarra} ITAM, versión de la carpeta --- 14 capítulos\par}
\vspace{0.2cm}
{\large\etiqueta{colorDP}{DP} \quad {\color{tonoOscuro}Thomas J. Sargent y John Stachurski}\par}
{\normalsize \emph{\color{colorDP}Dynamic Programming, Volume I: Finite States}\par}
{\small\color{tonoPizarra} Versión de abril de 2025 --- 10 capítulos\par}
\vspace{0.2cm}
{\large\etiqueta{colorMS}{MS} \quad {\color{tonoOscuro}Gjerrit Meinsma y Arjan van der Schaft}\par}
{\normalsize \emph{\color{colorMS}A Course on Optimal Control}\par}
{\small\color{tonoPizarra} Springer, 2023 --- 5 capítulos y 2 apéndices\par}
\vspace{0.5cm}
{\normalsize\color{colorPUCP} Con temas y lecturas de siete programas de la PUCP\par}
\vfill
{\setlength{\fboxsep}{6pt}\colorbox{tonoPrincipal}{\normalsize\sffamily\bfseries\color{white} Catorce bloques temáticos --- 8 meses}\par}
\vspace{0.2cm}
\end{titlepage}

\tableofcontents
\newpage

%------------------------------------------------------------------------------
\section*{Cómo usar este sílabo}
\addcontentsline{toc}{section}{Cómo usar este sílabo}

El plan está organizado en \textbf{catorce bloques temáticos}, agrupados en
dos partes de peso parecido. La \textbf{Parte I} (Bloques 1 a 8) construye la
matemática estática: lenguaje y orden, álgebra lineal, análisis en espacios
métricos, puntos fijos, cálculo diferencial, convexidad y optimización con
restricciones. La \textbf{Parte II} (Bloques 9 a 14) se dedica a la dinámica:
ecuaciones diferenciales y en diferencias, cálculo de variaciones, control
óptimo y programación dinámica, y cierra con aplicaciones y temas avanzados.
Conviene cerrar cada bloque antes de abrir el siguiente.

\medskip
\noindent\destacado{Por qué un curso aparte.} Los sílabos de microeconomía y de
macroeconomía no tienen pretramo matemático: suponen directamente este plan,
que estudia estas herramientas a fondo. La Parte II desarrolla en particular el control óptimo, el
cálculo de variaciones y la programación dinámica, que el sílabo de
macroeconomía usa desde su Bloque 1 y lleva a nivel de posgrado en su
Bloque 14, que se estudia antes de los modelos DSGE de su Bloque 12.

\medskip
\noindent\destacado{Relación con los demás sílabos.} Lo ideal es hacer este plan
antes que los de microeconomía y macroeconomía, o al menos la Parte I antes
del primero y la Parte II antes del segundo: los Bloques 9 y 10 antes de su
tema 1.6, el Bloque 12 antes de su Bloque 4 y el Bloque 13 antes de su
Bloque 14. El Bloque 14 del
sílabo de macroeconomía da por conocido el análisis real del capítulo~3 de
Stokey, Lucas y Prescott: los Bloques 3, 4 y 13 de este plan lo cubren. La
teoría de juegos queda fuera porque ocupa los Bloques 6 y 7 del sílabo de
microeconomía, y los métodos numéricos, porque corresponden a otro curso
(Anexo B).

\medskip
\noindent\destacado{Criterio general de reparto} (en el orden en que se lee cada tema):
\begin{itemize}[leftmargin=1.4em,itemsep=2pt]
  \item \textcolor{colorCG}{\textbf{Chávez y Gallardo}} es el texto base de la
        Parte I: es el libro del curso MAT291 de la PUCP, que dictan sus
        autores, y cubre el álgebra lineal, la topología, la convexidad y la
        optimización estática. Se lee \emph{primero}.
  \item \textcolor{colorCT}{\textbf{Carter}} es la profundización rigurosa de
        la Parte I y la lectura principal de lo que Chávez y Gallardo no
        tratan: órdenes, correspondencias, teoremas del máximo y de punto fijo,
        teoremas de separación en general y estática comparativa monótona.
  \item \textcolor{colorDF}{\textbf{De la Fuente}} aporta el cálculo
        diferencial y la función implícita en la Parte I, y es el texto base de
        los sistemas dinámicos y de las aplicaciones de la Parte II.
  \item \textcolor{colorKS}{\textbf{Kamien y Schwartz}} es la referencia del
        cálculo de variaciones y del control óptimo (Bloques 11 y 12);
        \textcolor{colorLR}{\textbf{Lomelí y Rumbos}} y
        \textcolor{colorBL}{\textbf{Bonifaz y Lama}} dan la versión en
        castellano, con aplicaciones.
  \item \textcolor{colorDP}{\textbf{Sargent y Stachurski}} es el texto
        principal de la programación dinámica moderna (Bloques 13 y 14);
        \textcolor{colorST}{\textbf{Stachurski}} lo extiende a espacios de
        estado generales.
  \item \textcolor{colorMS}{\textbf{Meinsma y van der Schaft}} conecta el
        cálculo de variaciones, el principio del máximo y la ecuación de
        Hamilton--Jacobi--Bellman, y cubre el control lineal-cuadrático.
\end{itemize}

\medskip
\noindent\destacado{Qué entra en las horas.} Las líneas \emph{Lecturas} y
\emph{Profundización} entran en el cómputo de horas. Las líneas
\emph{Consulta} y \emph{Consulta (PUCP)} remiten a libros de la carpeta y a
referencias de los programas que no entran (Anexo E).

\medskip
\noindent\destacado{Sobre los ejercicios.} La matemática se aprende demostrando y
resolviendo. De la Fuente trae más de doscientos problemas resueltos (su
apéndice), Kamien y Schwartz cierran cada sección con ejercicios, y Sargent y
Stachurski incluyen soluciones seleccionadas (su apéndice~C). Se recomienda
reservar al menos tres de las diez horas semanales para ejercicios.

\medskip
\noindent\destacado{Notación.} \CG{n.m} remite a la sección $m$ del capítulo $n$
de Chávez y Gallardo, según el índice del ejemplar físico; \CT{n.m}, \DF{n.m},
\LR{n.m}, \MS{n.m}, \DP{n.m}, \ST{n.m}, \LU{n.m} y \EL{n.m}, a las secciones de
los demás libros. De la Fuente subdivide sus secciones con letras: \DF{7.1c}
es el apartado c de la sección 1 del capítulo 7. Kamien y Schwartz numeran
las secciones dentro de cada parte: \KS{I.9} es la sección 9 de la Parte I
(cálculo de variaciones) y \KS{II.3}, la sección 3 de la Parte II (control
óptimo). Bonifaz y Lama numeran sus capítulos en romanos: \BL{III.2} es la
sección 2 del capítulo III.

%------------------------------------------------------------------------------
\newpage
\section*{Cronograma por semanas}
\addcontentsline{toc}{section}{Cronograma por semanas}

Calendario del plan: 340 horas distribuidas en \textbf{34 semanas de
contenido}, que equivalen a \textbf{8 meses de calendario}, a un ritmo de diez
horas semanales.

\medskip
\noindent Todas las semanas son de materia nueva y no hay semanas de colchón.
La columna de horas es la que gobierna el reparto. Si hay semanas perdidas por
feriados, viajes o imprevistos, el ajuste correcto es alargar el calendario,
no comprimir los bloques.

\medskip
{\filasalternas
\noindent\begin{tabularx}{\textwidth}{lXcc}
\toprule
\rowcolor{tonoOscuro}\cab{Semanas} & \cab{Contenido} & \cab{Sem.} & \cab{Horas} \\
\midrule
\rowcolor{tonoA}\multicolumn{4}{l}{\cab{\textit{Parte I --- Fundamentos y optimización estática}}} \\
1     & Bloque 1 --- Lenguaje matemático, conjuntos y orden & 1 & 10 \\
2--4  & Bloque 2 --- Álgebra lineal & 3 & 30 \\
5--6  & Bloque 3 --- Análisis en espacios métricos y normados & 2 & 20 \\
7--8  & Bloque 4 --- Correspondencias, teorema del máximo y puntos fijos & 2 & 20 \\
9--10 & Bloque 5 --- Cálculo diferencial en $\mathbb{R}^n$ & 2 & 20 \\
11--12 & Bloque 6 --- Convexidad & 2 & 20 \\
\rowcolor{tonoHielo}\textbf{13--15} & \textbf{Bloque 7 --- Optimización estática} & \textbf{3} & \textbf{30} \\
16    & Bloque 8 --- Estática comparativa y aplicaciones microeconómicas & 1 & 10 \\
\midrule
\rowcolor{tonoA}\multicolumn{4}{l}{\cab{\textit{Parte II --- Dinámica y optimización dinámica}}} \\
17--18 & Bloque 9 --- Ecuaciones diferenciales y en diferencias & 2 & 20 \\
19--20 & Bloque 10 --- Sistemas dinámicos multidimensionales & 2 & 20 \\
21--23 & Bloque 11 --- Cálculo de variaciones & 3 & 30 \\
\rowcolor{tonoHielo}\textbf{24--27} & \textbf{Bloque 12 --- Control óptimo} & \textbf{4} & \textbf{40} \\
\rowcolor{tonoHielo}\textbf{28--31} & \textbf{Bloque 13 --- Programación dinámica} & \textbf{4} & \textbf{40} \\
32--34 & Bloque 14 --- Aplicaciones y temas avanzados & 3 & 30 \\
\midrule
\rowcolor{tonoHielo}\multicolumn{2}{l}{\textbf{Total}} & \textbf{34} & \textbf{340} \\
\bottomrule
\end{tabularx}
}

\medskip
\noindent\destacado{Ritmo supuesto.} Diez horas por semana, todas de contenido
nuevo: por ejemplo, cuatro sesiones de dos horas entre semana para teoría y
una sesión de dos horas el sábado para ejercicios.

\medskip
\noindent\destacado{El reparto entre las dos partes.} La Parte I ocupa 16 semanas
y la Parte II, 18. El peso de la segunda es deliberado: el control óptimo y la
programación dinámica (Bloques 12 y 13, ocho semanas) son el núcleo del plan.

\medskip
\noindent\destacado{Los tramos difíciles.} Los Bloques 3 y 4 (semanas 5--8) son el
primer salto de abstracción: espacios métricos, compacidad y teoremas de punto
fijo. El segundo es el paso del control óptimo a la programación dinámica
(semanas 27--29), donde el teorema de la aplicación contractiva del tema 3.4
reaparece como herramienta central.

\medskip
\noindent\destacado{Hitos de control.} Conviene verificar el avance al terminar las
semanas 8, 16 y 27. Si en alguno de esos cortes se acumula más de una semana de
retraso, se alarga el calendario y no se comprime el bloque siguiente.

\subsection*{Detalle semana a semana}

{\small\filasalternas
\renewcommand{\arraystretch}{1.1}\noindent\begin{tabularx}{\textwidth}{clXc}
\toprule
\rowcolor{tonoOscuro}\cab{Sem.} & \cab{Bloque} & \cab{Temas} & \cab{Horas} \\
\midrule
1  & 1 & 1.1 Lógica, conjuntos, funciones y números reales; 1.2 Relaciones, órdenes y retículos & 10 \\
\midrule
2  & 2 & 2.1 Espacios y subespacios vectoriales, bases y dimensión; 2.2 Transformaciones lineales & 10 \\
3  & 2 & 2.3 Rango y sistemas de ecuaciones lineales; 2.4 Valores y vectores propios & 10 \\
4  & 2 & 2.5 Matrices simétricas y formas cuadráticas & 10 \\
\midrule
5  & 3 & 3.1 Normas, distancias y aproximación; 3.2 Sucesiones y topología de $\mathbb{R}^n$; 3.3 Continuidad & 10 \\
6  & 3 & 3.4 Completitud y teorema de la aplicación contractiva; 3.5 Compacidad, conexidad y teorema de Weierstrass & 10 \\
\midrule
7  & 4 & 4.1 Correspondencias y su continuidad; 4.2 El teorema del máximo & 10 \\
8  & 4 & 4.3 Teoremas de punto fijo; 4.4 Funciones monótonas y supermodularidad & 10 \\
\midrule
9  & 5 & 5.1 Derivada, derivadas parciales y jacobiano; 5.2 Taylor y teorema de la función inversa & 10 \\
10 & 5 & 5.3 Teorema de la función implícita y estática comparativa; 5.4 Funciones homogéneas & 10 \\
\midrule
11 & 6 & 6.1 Conjuntos convexos y conos; 6.2 Separación y lema de Farkas & 10 \\
12 & 6 & 6.3 Funciones cóncavas y convexas; 6.4 Funciones cuasicóncavas y cuasiconvexas & 10 \\
\midrule
13 & 7 & 7.1 El problema general y el enfoque geométrico; 7.2 Optimización sin restricciones & 10 \\
14 & 7 & 7.3 El problema de Lagrange; 7.4 El problema de Karush--Kuhn--Tucker: condiciones necesarias & 10 \\
15 & 7 & 7.5 Condiciones suficientes, programación cóncava y lineal; 7.6 Dualidad y punto silla & 10 \\
\midrule
16 & 8 & 8.1 Teorema de la envolvente y estática comparativa; 8.2 Preferencias, utilidad y equilibrio general & 10 \\
\bottomrule
\end{tabularx}

\medskip
{\filasalternas
\renewcommand{\arraystretch}{1.1}\noindent\begin{tabularx}{\textwidth}{clXc}
\toprule
\rowcolor{tonoOscuro}\cab{Sem.} & \cab{Bloque} & \cab{Temas} & \cab{Horas} \\
\midrule
17 & 9 & 9.1 Conceptos básicos, existencia y unicidad; 9.2 Ecuaciones diferenciales escalares & 10 \\
18 & 9 & 9.3 Ecuaciones en diferencias escalares; 9.4 Ecuaciones en diferencias estocásticas & 10 \\
\midrule
19 & 10 & 10.1 Sistemas lineales continuos y discretos; 10.2 Sistemas no lineales en el plano & 10 \\
20 & 10 & 10.3 Estabilidad y funciones de Liapunov; 10.4 Aplicaciones: IS-LM, previsión perfecta, Solow y OLG; 10.5 Transformada de Laplace & 10 \\
\midrule
21 & 11 & 11.1 La ecuación de Euler--Lagrange; 11.2 Segundo orden y restricciones integrales & 10 \\
22 & 11 & 11.3 Transversalidad y extremos libres; 11.4 Horizonte infinito y diagramas de fase & 10 \\
23 & 11 & 11.5 Esquinas y restricciones de desigualdad; 11.6 Optimización en espacios de funciones & 10 \\
\midrule
24 & 12 & 12.1 El principio del máximo; 12.2 Condiciones suficientes & 10 \\
25 & 12 & 12.3 Condiciones terminales y varias variables; 12.4 Descuento y horizonte infinito & 10 \\
26 & 12 & 12.5 Restricciones sobre controles y estados; 12.6 Saltos, existencia y formulaciones alternativas & 10 \\
27 & 12 & 12.7 Control lineal-cuadrático y ecuación de Riccati; 12.8 Interpretación económica y aplicaciones & 10 \\
\midrule
28 & 13 & 13.1 Principio de optimalidad y ecuación de Bellman; 13.2 Horizonte infinito y contracciones & 10 \\
29 & 13 & 13.3 Ecuación de Euler y métodos de solución; 13.4 Operadores, orden y puntos fijos & 10 \\
30 & 13 & 13.5 Cadenas de Markov y procesos de decisión de Markov; 13.6 Parada óptima & 10 \\
31 & 13 & 13.7 La ecuación de Hamilton--Jacobi--Bellman; 13.8 Espacios de estado generales & 10 \\
\midrule
32 & 14 & 14.1 Búsqueda de empleo; 14.2 Crecimiento óptimo en tiempo discreto y continuo & 10 \\
33 & 14 & 14.3 Inversión con costos de ajuste; 14.4 Descuento estocástico y preferencias recursivas & 10 \\
34 & 14 & 14.5 Procesos de decisión recursivos, robustez y ambigüedad; 14.6 Decisión en tiempo continuo & 10 \\
\bottomrule
\end{tabularx}
}
}

%==============================================================================
\parte{Parte I --- Fundamentos y optimización estática}

%------------------------------------------------------------------------------
\bloque{Bloque 1 --- Lenguaje matemático, conjuntos y orden --- 1 semana}
{El lenguaje con el que está escrito todo lo demás: lógica, métodos de
demostración, conjuntos y relaciones. Las relaciones de orden reaparecen en
las preferencias (tema 8.2), en la estática comparativa monótona (tema 4.4) y
en la programación dinámica (tema 13.4). Cubre el primer tema de MAT659.}

\tema{1.1 Lógica, conjuntos, funciones y números reales}
\lecturas{\DF{1.1--1.2} \quad \DF{1.4} \quad \DF{1.6}}
\apoyo{\CT{1.1}}
\consultalib{Corbae, Stinchcombe y Zeman (2009), caps.~1--3}
\nota{Cuantificadores, implicación y métodos de prueba (directa, por
contradicción, por inducción). La propiedad del supremo es la base de la
completitud del tema 3.4.}

\tema{1.2 Relaciones, órdenes y retículos}
\lecturas{\CT{1.2} \quad \DF{1.3}}
\consulta{Bridges y Mehta (1995)}
\nota{Relaciones de equivalencia, órdenes parciales y débiles, retículos y el
orden de Pareto. \CT{1.2.6} adelanta el problema de la agregación de
preferencias.}

%------------------------------------------------------------------------------
\bloque{Bloque 2 --- Álgebra lineal --- 3 semanas}
{Capítulos 1 a 3 de Chávez y Gallardo, con De la Fuente y Carter como
profundización. Cubre el Módulo 1 de 1ECO75 (independencia lineal,
sistemas, diagonalización y formas cuadráticas).}

\tema{2.1 Espacios y subespacios vectoriales, bases y dimensión}
\lecturas{\CG{1.1--1.3}}
\apoyo{\DF{3.1} \quad \CT{1.4.1--1.4.2}}
\consultalib{Simon y Blume (1994), caps.~10--11 y 27}

\tema{2.2 Transformaciones lineales}
\lecturas{\CG{1.4}}
\apoyo{\DF{3.2--3.4} \quad \CT{3.1--3.3}}
\nota{Núcleo, imagen e isomorfismos; las funciones lineales entre espacios
normados y su norma (\DF{3.4}) se usan en el tema 5.1. \CT{3.3.2} introduce
los hiperplanos del tema 6.2.}

\tema{2.3 Rango y sistemas de ecuaciones lineales}
\lecturas{\CG{2.1--2.3}}
\apoyo{\CT{3.6}}
\consultalib{Simon y Blume (1994), caps.~7--9}
\nota{\CT{3.6.2} trata los sistemas de desigualdades lineales, preparación del
lema de Farkas (tema 6.2); \CT{3.6.3--3.6.4} aplica la teoría a los modelos de
insumo-producto y a las cadenas de Markov, que vuelven en el tema 13.5.}

\tema{2.4 Valores y vectores propios}
\lecturas{\CG{3.1--3.2}}
\apoyo{\DF{3.5--3.6} \quad \CT{3.5.1--3.5.2}}
\consultalib{Simon y Blume (1994), cap.~23}
\nota{Semejanza de matrices y diagonalización. Es la herramienta con la que se
resuelven los sistemas dinámicos lineales del tema 10.1.}

\tema{2.5 Matrices simétricas y formas cuadráticas}
\lecturas{\CG{3.3}}
\apoyo{\CT{3.4} \quad \CT{3.5.3} \quad \DF{6}, apéndice (formas cuadráticas)}
\consultalib{Simon y Blume (1994), cap.~16}
\nota{Formas bilineales, productos internos y criterios de definición de una
matriz. Son las condiciones de segundo orden de los Bloques 7 y 11.}

%------------------------------------------------------------------------------
\bloque{Bloque 3 --- Análisis en espacios métricos y normados --- 2 semanas}
{El capítulo 4 de Chávez y Gallardo en $\mathbb{R}^n$ y su generalización a
espacios métricos con De la Fuente y Carter. Cubre el Módulo 2 de 1ECO75.
Es el análisis real que el sílabo de macroeconomía (tema 14.1) da por
conocido al leer a Stokey, Lucas y Prescott.}

\tema{3.1 Normas, distancias y aproximación}
\lecturas{\CG{4.1--4.2} \quad \DF{2.1}}
\apoyo{\CT{1.5}}
\nota{Espacios métricos y normados; las normas equivalentes de \DF{2.10} se
leen en el tema 3.5.}

\tema{3.2 Sucesiones y topología de $\mathbb{R}^n$}
\lecturas{\CG{4.3} \quad \DF{2.2--2.4}}
\apoyo{\CT{1.3.1}}
\nota{Convergencia, conjuntos abiertos y cerrados, interior, clausura y
frontera.}

\tema{3.3 Continuidad}
\lecturas{\DF{2.5--2.6}}
\apoyo{\CT{2.3.1--2.3.3}}
\nota{Límites, continuidad, semicontinuidad y continuidad uniforme. El
vocabulario de funciones de \CT{2.1.1--2.1.4} se consulta según haga falta.}

\tema{3.4 Completitud y teorema de la aplicación contractiva}
\lecturas{\DF{2.7}}
\apoyo{\CT{1.3.2} \quad \ST{3.1--3.2}}
\consultalib{Acemoglu (2009), secc.~6.3}
\nota{Sucesiones de Cauchy, espacios completos y el teorema de Banach. Es la
base de la programación dinámica en horizonte infinito (tema 13.2).}

\tema{3.5 Compacidad, conexidad y teorema de Weierstrass}
\lecturas{\DF{2.8--2.10}}
\consultalib{Corbae, Stinchcombe y Zeman (2009), caps.~4 y 6 \quad Acemoglu (2009), apéndice A.1--A.4}
\nota{Caracterizaciones de la compacidad y existencia de máximos de funciones
continuas sobre compactos: el resultado de existencia del tema 7.1.}

%------------------------------------------------------------------------------
\bloque{Bloque 4 --- Correspondencias, teorema del máximo y puntos fijos --- 2 semanas}
{Los «tres teoremas fundamentales de la economía matemática» que Carter
destaca: el teorema del máximo, el de punto fijo de Brouwer y el del
hiperplano separador (este último, en el tema 6.2). Aquí Carter es lectura
principal. Cubre el segundo y el tercer tema de MAT659.}

\tema{4.1 Correspondencias y su continuidad}
\lecturas{\CT{2.1.5} \quad \CT{2.3.4} \quad \DF{2.11}}
\consultalib{Acemoglu (2009), apéndice A.5}
\nota{Hemicontinuidad superior e inferior y gráfica cerrada. Las
correspondencias de demanda y de mejor respuesta son el ejemplo económico.}

\tema{4.2 El teorema del máximo}
\lecturas{\CT{2.3.5} \quad \DF{7.2a}}
\consulta{Rockafellar y Wets (2009) \quad Royset (2020)}
\nota{Continuidad de la función valor y hemicontinuidad superior de la
correspondencia de maximizadores (Berge). La generalización a convergencia
de conjuntos (hipertopologías) de MAT659 queda en el Anexo D.3.}

\tema{4.3 Teoremas de punto fijo}
\lecturas{\CT{2.4} \quad \DF{5.3}}
\nota{Tarski, Banach y Brouwer, con Kakutani como extensión a
correspondencias. Son los resultados de existencia del equilibrio general
(sílabo de microeconomía, tema 9.3) y del equilibrio de Nash (tema 6.4 de ese
sílabo).}

\tema{4.4 Funciones monótonas, supermodularidad y teorema del máximo monótono}
\lecturas{\CT{2.2}}
\apoyo{\DP{2.2}}
\consulta{Topkis (1998), caps.~2--3}
\nota{Funciones supermodulares en retículos y soluciones óptimas crecientes:
la estática comparativa sin derivadas que MAT659 desarrolla en su tercera
parte. Se aplica en el tema 8.1. \DP{2.2} trata el orden puntual y la
monotonía paramétrica que reaparecen en la programación dinámica.}

%------------------------------------------------------------------------------
\bloque{Bloque 5 --- Cálculo diferencial en $\mathbb{R}^n$ --- 2 semanas}
{Chávez y Gallardo no dedican un capítulo al cálculo; aquí De la Fuente es el
texto principal y Carter la profundización. Cubre la Sesión 2.4 de 1ECO75 y
las series de Taylor de su sumilla.}

\tema{5.1 Derivada, derivadas parciales y jacobiano}
\lecturas{\DF{4.1--4.3}}
\apoyo{\CT{4.1--4.3}}
\consultalib{Simon y Blume (1994), caps.~14 y 30}
\nota{La derivada como aproximación lineal (Fréchet), derivadas direccionales
y teorema del valor medio.}

\tema{5.2 Diferenciabilidad continua, Taylor y teorema de la función inversa}
\lecturas{\DF{4.4}}
\apoyo{\CT{4.4} \quad \CT{4.5.1}}
\consultalib{Bonifaz y Winkelried (2003), cap.~III}

\tema{5.3 Teorema de la función implícita y estática comparativa}
\lecturas{\DF{5.1--5.2}}
\apoyo{\CT{4.5.2} \quad \CT{6.3}}
\consultalib{Simon y Blume (1994), cap.~15}
\nota{Modelos lineales, derivadas de funciones implícitas y estática
comparativa de modelos de equilibrio.}

\tema{5.4 Funciones homogéneas y homotéticas}
\lecturas{\DF{4.5}}
\apoyo{\CT{3.8} \quad \CT{4.6.2}}
\consultalib{Simon y Blume (1994), cap.~20}

%------------------------------------------------------------------------------
\bloque{Bloque 6 --- Convexidad --- 2 semanas}
{Capítulos 5 y 6 de Chávez y Gallardo. Cubre las primeras unidades de MAT693
(conjuntos convexos, separación, concavidad generalizada) y las dos primeras
semanas de MAT291.}

\tema{6.1 Conjuntos convexos y conos}
\lecturas{\CG{5.1} \quad \DF{6.1a--c}}
\apoyo{\CT{1.4.3--1.4.5}}
\consultalib{Boyd y Vandenberghe (2004), cap.~2}
\nota{Conjuntos afines, combinaciones y envolturas convexas, conos convexos e
interior relativo.}

\tema{6.2 Separación de conjuntos convexos y lema de Farkas}
\lecturas{\CG{5.2--5.3} \quad \DF{6.1d}}
\apoyo{\CT{3.9} (hasta 3.9.3)}
\nota{Hiperplanos separadores y de soporte, el teorema de Hahn--Banach y los
teoremas de la alternativa. El lema de Farkas es la pieza que justifica las
condiciones de Kuhn--Tucker (tema 7.4).}

\tema{6.3 Funciones cóncavas y convexas}
\lecturas{\CG{6.1--6.2} \quad \DF{6.2}}
\apoyo{\CT{3.7.1} \quad \CT{4.6.1}}
\consultalib{Boyd y Vandenberghe (2004), cap.~3}
\nota{Caracterizaciones con el epígrafo, la derivada y el hessiano.}

\tema{6.4 Funciones cuasicóncavas y cuasiconvexas}
\lecturas{\CG{6.3} \quad \DF{6.3}}
\apoyo{\CT{3.7.2--3.7.3}}
\consultalib{Simon y Blume (1994), cap.~21}
\nota{La cuasiconcavidad es la traducción matemática de las preferencias
convexas. \CT{3.7.3} da los teoremas del máximo bajo convexidad.}

%------------------------------------------------------------------------------
\bloque{Bloque 7 --- Optimización estática --- 3 semanas}
{Capítulos 7 a 9 de Chávez y Gallardo, el núcleo de MAT291 (Lagrange y
Karush--Kuhn--Tucker), del Módulo 3 de 1ECO75 y de las unidades centrales de
MAT693.}

\tema{7.1 El problema general y el enfoque geométrico}
\lecturas{\CG{7.1--7.2}}
\apoyo{\CT{5.1}}
\nota{Existencia (Weierstrass, tema 3.5), máximos locales y globales y
optimización sobre un conjunto convexo (\DF{7.1a}).}

\tema{7.2 Optimización sin restricciones: soluciones interiores}
\lecturas{\CG{7.3}}
\apoyo{\CT{5.2}}
\consultalib{Simon y Blume (1994), cap.~17}
\nota{Condiciones de primer y segundo orden; las formas cuadráticas del tema
2.5 deciden la naturaleza del punto crítico.}

\tema{7.3 El problema de Lagrange}
\lecturas{\CG{8.1--8.2} \quad \DF{7.1b}}
\apoyo{\CT{5.3}}
\consultalib{Simon y Blume (1994), caps.~18--19}
\nota{Carter da cuatro enfoques del multiplicador (perturbación, geométrico,
función implícita y lagrangiano) y su lectura como precio sombra
(\CT{5.3.5}).}

\tema{7.4 El problema de Karush--Kuhn--Tucker: condiciones necesarias}
\lecturas{\CG{9.1} \quad \CG{9.3} \quad \DF{7.1c}}
\apoyo{\CT{5.4.1--5.4.2}}
\nota{El problema mixto, las condiciones de no negatividad y las
calificaciones de las restricciones.}

\tema{7.5 Condiciones suficientes, programación cóncava y programación lineal}
\lecturas{\CG{9.2} \quad \DF{7.1d}}
\apoyo{\CT{5.4.3--5.4.5}}
\nota{Suficiencia bajo concavidad y cuasiconcavidad; programación cóncava sin
diferenciabilidad.}

\tema{7.6 Dualidad lagrangiana y punto silla}
\lecturas{\CT{3.9.2} \quad \CT{3.7.4}}
\apoyo{\LU{8.1--8.6}}
\consultalib{Boyd y Vandenberghe (2004), cap.~5}
\nota{Teoremas minimax, punto silla del lagrangiano y dualidad en espacios
vectoriales (Luenberger). Es la unidad e de MAT693 y la puerta al
tratamiento funcional del tema 11.6.}

%------------------------------------------------------------------------------
\bloque{Bloque 8 --- Estática comparativa y aplicaciones microeconómicas --- 1 semana}
{El cierre de la Parte I: cómo cambia la solución de un problema de
optimización cuando cambian los parámetros, y la aplicación a la teoría del
consumidor. Capítulo 6 de Carter y capítulo 10 de Chávez y Gallardo.}

\tema{8.1 Teorema de la envolvente y estática comparativa}
\lecturas{\CG{8.3} \quad \CT{6.1--6.2} \quad \DF{7.2b--c}}
\nota{Tres enfoques: preferencia revelada, función valor (envolvente) y
monotonía (que aplica el tema 4.4). Es la unidad f de MAT693 y la
herramienta de la dualidad del Bloque 2 del sílabo de microeconomía.}

\tema{8.2 Preferencias, utilidad y equilibrio general}
\lecturas{\CG{10.1--10.3} \quad \DF{8.1--8.3}}
\apoyo{\CT{1.6}}
\nota{Representación de preferencias por funciones de utilidad, demanda y
existencia del equilibrio walrasiano con los teoremas de punto fijo del
tema 4.3. Se lee como aplicación de la matemática: el desarrollo económico
está en los Bloques 1, 2 y 9 del sílabo de microeconomía.}

%==============================================================================
\parte{Parte II --- Dinámica y optimización dinámica}

%------------------------------------------------------------------------------
\bloque{Bloque 9 --- Ecuaciones diferenciales y en diferencias --- 2 semanas}
{Primera parte de la dinámica de MAT291 (temas d y e) y de la sumilla de
1ECO75. De la Fuente da la teoría y Lomelí y Rumbos, la versión en castellano
con aplicaciones.}

\tema{9.1 Conceptos básicos, existencia y unicidad}
\lecturas{\DF{9.1--9.2} \quad \DF{9.6}}
\consultalib{Acemoglu (2009), apéndice B}
\nota{Problemas de valor inicial y de frontera, flujo de un sistema autónomo,
estados estacionarios y estabilidad; existencia local, soluciones maximales y
dependencia de las condiciones iniciales.}

\tema{9.2 Ecuaciones diferenciales escalares}
\lecturas{\DF{9.3} \quad \DF{9.5} \quad \LR{2--3}}
\consultalib{Bonifaz y Winkelried (2003), cap.~IV \quad Simon y Blume (1994), cap.~24}
\nota{Ecuaciones lineales de coeficientes constantes, ecuaciones no lineales
autónomas, diagramas de fase de una variable y dinámica comparativa.}

\tema{9.3 Ecuaciones en diferencias escalares}
\lecturas{\DF{9.4} \quad \LR{6}}
\apoyo{\EL{1.1--1.3} \quad \EL{1.5--1.7}}
\consultalib{Bonifaz y Winkelried (2003), cap.~VI}
\nota{Estabilidad de puntos fijos, ciclos, bifurcaciones y caos en el modelo
logístico.}

\tema{9.4 Ecuaciones en diferencias estocásticas y expectativas racionales}
\lecturas{\LR{9}}
\nota{Iteración hacia adelante, burbujas y forma reducida: la entrada a los
modelos con expectativas racionales del Bloque 9 del sílabo de
macroeconomía.}

%------------------------------------------------------------------------------
\bloque{Bloque 10 --- Sistemas dinámicos multidimensionales --- 2 semanas}
{Sistemas lineales y no lineales, diagramas de fase en el plano y
estabilidad. Cubre las Sesiones 4.1 y 4.2 de 1ECO75 y el tema de sistemas
dinámicos de MAT291.}

\tema{10.1 Sistemas lineales en tiempo continuo y discreto}
\lecturas{\DF{10.1--10.2} \quad \LR{4} \quad \LR{7}}
\apoyo{\EL{3.1--3.3}}
\consultalib{Simon y Blume (1994), cap.~25}
\nota{Solución por diagonalización (tema 2.4), valores propios complejos y
repetidos, espacios estable e inestable.}

\tema{10.2 Sistemas no lineales en el plano: diagramas de fase y linealización}
\lecturas{\DF{10.3} \quad \LR{5}}
\nota{Clasificación de los puntos de equilibrio de los sistemas $2\times 2$ y
análisis de puntos silla, la estructura de casi todos los modelos
macroeconómicos con previsión perfecta.}

\tema{10.3 Estabilidad y funciones de Liapunov}
\lecturas{\MS{B.1--B.4} \quad \LR{8}}
\apoyo{\EL{4.1--4.6}}
\nota{Definiciones de estabilidad, el método directo de Liapunov y el
principio de invariancia de LaSalle. \LR{8} lo aplica a la estabilidad del
equilibrio walrasiano (el tâtonnement).}

\tema{10.4 Aplicaciones: IS-LM dinámico, previsión perfecta, Solow y OLG}
\lecturas{\DF{11.1--11.4}}
\apoyo{\ST{4.1}}
\nota{Dornbusch, un modelo de precios de acciones, Solow y Diamond. \ST{4.1}
añade la estabilidad global y la equivalencia de dinámicas por conjugación.}

\tema{10.5 La transformada de Laplace y la dinámica de transición}
\lecturas{Heijdra (2017), apéndice matemático, secc.~A.6}
\nota{Definición y propiedades de la transformada de Laplace, y su uso para
resolver sistemas lineales de ecuaciones diferenciales con choques que cambian
en el tiempo. Es la técnica con que Heijdra obtiene la dinámica de transición
en el modelo de Ramsey (cap.~13), que el sílabo de macroeconomía lee en los
temas 4.7 y 4.8.}

%------------------------------------------------------------------------------
\bloque{Bloque 11 --- Cálculo de variaciones --- 3 semanas}
{El tema f de MAT291 y la sumilla de 1ECO75 (ecuación de Euler--Lagrange y
condiciones de transversalidad). De la Fuente no lo trata; el texto base es
la Parte I de Kamien y Schwartz, con Meinsma y van der Schaft para la
teoría y Lomelí y Rumbos y Bonifaz y Lama para las aplicaciones.}

\tema{11.1 El problema más simple y la ecuación de Euler--Lagrange}
\lecturas{\KS{I.1--I.5} \quad \LR{11.1--11.2}}
\apoyo{\MS{1.1--1.3} \quad \BL{I} \quad \BL{II.1}}
\consulta{Cerdá (2001) \quad Chiang (1992)}
\nota{Derivación por variaciones admisibles, casos especiales de la ecuación de
Euler, la identidad de Beltrami y la extracción óptima de recursos.}

\tema{11.2 Condiciones de segundo orden y restricciones integrales}
\lecturas{\KS{I.6--I.7} \quad \LR{11.5}}
\apoyo{\MS{1.6--1.7} \quad \BL{II.3}}
\nota{Condición de Legendre, suficiencia bajo concavidad y el problema
isoperimétrico.}

\tema{11.3 Condiciones de transversalidad y extremos libres}
\lecturas{\KS{I.8--I.12} \quad \LR{11.4} \quad \LR{11.6}}
\apoyo{\MS{1.4--1.5} \quad \BL{II.2}}
\nota{Valor final libre, horizonte libre, extremos restringidos, valor de
rescate y análisis de sensibilidad.}

\tema{11.4 Horizonte infinito y diagramas de fase}
\lecturas{\KS{I.15--I.17} \quad \LR{11.3} \quad \LR{11.7}}
\apoyo{\BL{II.4}}
\nota{Problemas autónomos, trayectorias de aproximación más rápida, análisis
diagramático y el modelo de Ramsey como problema variacional.}

\tema{11.5 Esquinas, restricciones de desigualdad y varias funciones}
\lecturas{\KS{I.13--I.14} \quad \KS{I.18}}
\nota{Condiciones de Weierstrass--Erdmann y restricciones en $(t,x)$: el
límite del cálculo de variaciones que motiva el control óptimo.}

\tema{11.6 El cálculo de variaciones como optimización en espacios de funciones}
\lecturas{\LU{7.1--7.6}}
\nota{Diferenciales de Gâteaux y de Fréchet, extremos de funcionales y la
ecuación de Euler--Lagrange como condición de primer orden en un espacio
normado. Une el Bloque 7 con el Bloque 11.}

%------------------------------------------------------------------------------
\bloque{Bloque 12 --- Control óptimo --- 4 semanas}
{El núcleo de la Parte II. Cubre el tema f de MAT291, las Sesiones 4.3--4.4 de
1ECO75 y el control determinístico de MAT835. Es el control óptimo que el
sílabo de macroeconomía supone desde su Bloque 4, donde se usa para el modelo
de Ramsey. Kamien y Schwartz (Parte II) es el
texto principal.}

\tema{12.1 El principio del máximo}
\lecturas{\DF{12.2a} \quad \KS{II.1--II.2} \quad \LR{12.1--12.2}}
\apoyo{\MS{2.1--2.5} \quad \BL{III.1}}
\consulta{Cerdá (2001) \quad Léonard y Van Long (1992)}
\nota{Hamiltoniano, variable de coestado y condiciones necesarias. \MS{2}
deriva el principio del mínimo desde los multiplicadores de Lagrange.}

\tema{12.2 Condiciones suficientes}
\lecturas{\DF{12.2b} \quad \KS{II.3} \quad \KS{II.15}}
\apoyo{\BL{III.3} \quad \MS{2.8}}
\nota{Teoremas de suficiencia de Mangasarian y de Arrow, y el papel de la
convexidad (Bloque 6).}

\tema{12.3 Condiciones terminales y varias variables}
\lecturas{\KS{II.5--II.7} \quad \LR{12.3} \quad \LR{12.6}}
\apoyo{\MS{2.6--2.7} \quad \BL{III.2}}
\nota{Punto final fijo, libre o restringido; horizonte libre; varios estados y
controles.}

\tema{12.4 Descuento, hamiltoniano a valor corriente y horizonte infinito}
\lecturas{\KS{II.8--II.9} \quad \LR{12.4--12.5}}
\apoyo{\BL{III.5--III.6}}
\consultalib{Acemoglu (2009), secc.~7.3--7.5}
\nota{Dinámica comparativa, equilibrios de problemas autónomos y la
controversia sobre la condición de transversalidad en horizonte infinito.}

\tema{12.5 Restricciones sobre controles y estados}
\lecturas{\DF{12.2c} \quad \KS{II.10--II.13} \quad \KS{II.17} \quad \LR{13}}
\nota{Controles acotados, soluciones \emph{bang-bang} y singulares,
restricciones mixtas e integrales y restricciones de desigualdad sobre el
estado.}

\tema{12.6 Saltos, existencia y formulaciones alternativas}
\lecturas{\KS{II.14} \quad \KS{II.16} \quad \KS{II.18}}
\apoyo{\LU{9.5--9.6}}
\nota{El principio de Pontryagin en su forma general y la existencia de
soluciones; saltos en el estado y cambios de régimen. Luenberger lo deduce
de la teoría local de la optimización con restricciones (tema 7.6).}

\tema{12.7 Control lineal-cuadrático y ecuación de Riccati}
\lecturas{\MS{4.1--4.5}}
\nota{El regulador lineal-cuadrático por el principio del mínimo y por
programación dinámica, la ecuación diferencial de Riccati y su versión
algebraica en horizonte infinito. Prepara el tema 14.2 del sílabo de
macroeconomía.}

\tema{12.8 Interpretación económica y aplicaciones}
\lecturas{\LR{12.7--12.8} \quad \KS{II.4} \quad \BL{III.4}}
\consultalib{Acemoglu (2009), secc.~7.7--7.8}
\nota{El coestado como precio sombra (Dorfman), la comparación con el cálculo
de variaciones, la contaminación óptima y el modelo neoclásico de
crecimiento (\BL{III.6}).}

%------------------------------------------------------------------------------
\bloque{Bloque 13 --- Programación dinámica --- 4 semanas}
{El tema g de MAT291 y la programación dinámica de 1ECO75, MAT616 y MAT835.
Sargent y Stachurski es el texto principal: trabaja con estados finitos, lo
que permite demostrar todo con álgebra lineal y puntos fijos (Bloques 2 a 4).
El sílabo de macroeconomía (Bloque 14) retoma el tema con Stokey, Lucas y
Prescott y lo lleva al cómputo antes de estudiar los modelos DSGE.}

\tema{13.1 Principio de optimalidad y ecuación de Bellman}
\lecturas{\DF{12.1a} \quad \DP{1.1} \quad \LR{14.1--14.3}}
\apoyo{\MS{3.1--3.3} \quad \BL{IV.1--IV.2}}
\nota{Horizonte finito e inducción hacia atrás; el problema de búsqueda de
empleo como ejemplo guía.}

\tema{13.2 Horizonte infinito, operador de Bellman y contracciones}
\lecturas{\DF{12.1b} \quad \DP{1.2--1.3} \quad \LR{14.4}}
\consultalib{Acemoglu (2009), secc.~6.1--6.5}
\consulta{Stokey, Lucas y Prescott (1989), caps.~3--4}
\nota{Problemas estacionarios descontados: existencia y unicidad de la función
valor por el teorema de la aplicación contractiva (tema 3.4) y política
óptima.}

\tema{13.3 Ecuación de Euler y métodos de solución}
\lecturas{\BL{IV.2--IV.3} \quad \LR{14.5--14.6}}
\nota{La condición de Benveniste--Scheinkman, la ecuación de Euler y la
transversalidad; aproximaciones sucesivas y el método de «adivinar y
verificar» en el modelo de Ramsey discreto.}

\tema{13.4 Operadores, orden y puntos fijos}
\lecturas{\DP{2.1} \quad \DP{2.3}}
\nota{Estabilidad local, tasas de convergencia, matrices no negativas y
operadores lineales; junto con el orden de \DP{2.2} (tema 4.4), es la caja de
herramientas del resto del libro.}

\tema{13.5 Cadenas de Markov y procesos de decisión de Markov}
\lecturas{\DP{3} \quad \DP{5.1--5.2}}
\apoyo{\ST{4.2--4.3} \quad \LR{14.7}}
\nota{Distribuciones estacionarias y ergodicidad, esperanzas condicionales y el
modelo general de decisión de Markov: optimalidad, iteración de la función
valor y de la política, ahorro e inversión óptimos.}

\tema{13.6 Parada óptima}
\lecturas{\DP{4}}
\consulta{Øksendal (2010)}
\nota{Valoración de una empresa con salida, opciones americanas e
investigación y desarrollo. Es la versión en tiempo discreto de la primera
unidad de MAT835.}

\tema{13.7 La ecuación de Hamilton--Jacobi--Bellman}
\lecturas{\MS{3.4--3.6} \quad \KS{II.21}}
\nota{La programación dinámica en tiempo continuo, su relación con el
principio del máximo (el coestado como derivada de la función valor) y las
funciones de Liapunov de costo restante. Es la segunda unidad de MAT835.}

\tema{13.8 Programación dinámica en espacios de estado generales}
\lecturas{\ST{10.1} \quad \ST{12.1}}
\consultalib{Acemoglu (2009), cap.~16}
\nota{El problema con estados continuos y choques, optimalidad, y
monotonía y concavidad de la función valor y de la política. Las nociones de
teoría de la medida que usa se consultan en \ST{7} y \ST{9} (Anexo D.1).}

%------------------------------------------------------------------------------
\bloque{Bloque 14 --- Aplicaciones y temas avanzados --- 3 semanas}
{Primero, las aplicaciones clásicas de De la Fuente, resueltas con las tres
técnicas de la Parte II; después, los temas avanzados de Sargent y Stachurski:
descuento estocástico, preferencias recursivas, robustez y tiempo continuo.}

\tema{14.1 Búsqueda de empleo}
\lecturas{\DF{13.1}}
\apoyo{\DP{3.3}}
\nota{El modelo básico de búsqueda y un modelo macroeconómico basado en
búsqueda; \DP{3.3} añade el estado markoviano y la separación.}

\tema{14.2 Crecimiento óptimo en tiempo discreto y continuo}
\lecturas{\DF{13.2} \quad \DF{13.4}}
\consultalib{Acemoglu (2009), secc.~6.6 y 7.7}
\nota{La función de política y la ecuación de Euler del crecimiento óptimo
discreto; el modelo de Cass--Koopmans con impuestos y el costo de bienestar
de la tributación de factores.}

\tema{14.3 Inversión con costos de ajuste}
\lecturas{\DF{13.3}}
\apoyo{\DP{5.2.3}}
\consultalib{Acemoglu (2009), secc.~7.8}
\nota{La $q$ de Tobin y la acumulación de capital en una economía pequeña y
abierta. Enlaza con el tema 5.4 del sílabo de economía financiera.}

\tema{14.4 Descuento estocástico y preferencias recursivas}
\lecturas{\DP{6.1--6.3} \quad \DP{7.1--7.2}}
\nota{Factores de descuento variables y la condición del radio espectral,
valoración de activos, el teorema de Knaster--Tarski en espacios de
funciones y las preferencias de Epstein--Zin y sensibles al riesgo. Es la
base matemática del tema 5.3 del sílabo de economía financiera.}

\tema{14.5 Procesos de decisión recursivos, robustez y ambigüedad}
\lecturas{\DP{7.3} \quad \DP{8}}
\apoyo{\DP{9}}
\nota{Operadores de Koopmans, condiciones del tipo de Blackwell, procesos
contractivos y eventualmente contractivos, agentes adversarios, robustez y
ambigüedad; \DP{9} los unifica en la programación dinámica abstracta.}

\tema{14.6 Decisión en tiempo continuo}
\lecturas{\DP{10}}
\consulta{Fleming y Soner (2010) \quad Pham (2010)}
\nota{Cadenas de Markov en tiempo continuo, semigrupos y procesos de decisión
en tiempo continuo, con el problema de búsqueda como aplicación. Es el
puente hacia el control óptimo estocástico de MAT835 (Anexo D.2).}

%------------------------------------------------------------------------------
\newpage
\section*{Anexo A. De los capítulos a los bloques}
\addcontentsline{toc}{section}{Anexo A. De los capítulos a los bloques}

Capítulos de cada libro que entran en el plan y temas donde se leen. Lo que
queda fuera se detalla en el Anexo B.

\newcommand{\mapa}[2]{%
\subsection*{#1}
{\small\filasalternas
\noindent\begin{tabularx}{\textwidth}{>{\raggedright\arraybackslash}p{1.3cm}X>{\raggedright\arraybackslash}p{3.6cm}}
\toprule
\rowcolor{tonoOscuro}\cab{Cap.} & \cab{Título} & \cab{Temas} \\
\midrule
#2
\bottomrule
\end{tabularx}
}}

\mapa{A.1 Chávez y Gallardo, \emph{Álgebra lineal y optimización para el análisis económico} (texto base de la Parte I)}{%
1 & Elementos de álgebra lineal & 2.1, 2.2 \\
2 & Sistema lineal de ecuaciones algebraicas & 2.3 \\
3 & Vectores y valores propios & 2.4, 2.5 \\
4 & Espacios vectoriales topológicos & 3.1, 3.2 \\
5 & Conjuntos convexos & 6.1, 6.2 \\
6 & Funciones convexas y cóncavas & 6.3, 6.4 \\
7 & Introducción a la optimización & 7.1, 7.2 \\
8 & El problema de Lagrange & 7.3, 8.1 \\
9 & El problema de Karush--Kuhn--Tucker & 7.4, 7.5 \\
10 & Relaciones de preferencia y equilibrio general & 8.2 \\
}

\mapa{A.2 Carter, \emph{Foundations of Mathematical Economics}}{%
1 & Sets and Spaces (1.1--1.6) & 1.1, 1.2, 2.1, 3.1, 3.2, 3.4, 6.1, 8.2 \\
2 & Functions (2.1--2.4) & 3.3, 4.1--4.4 \\
3 & Linear Functions (3.1--3.9) & 2.2--2.5, 5.4, 6.2--6.4, 7.6 \\
4 & Smooth Functions (4.1--4.6) & 5.1--5.4, 6.3 \\
5 & Optimization (5.1--5.4) & 7.1--7.5 \\
6 & Comparative Statics (6.1--6.3) & 5.3, 8.1 \\
}

\mapa{A.3 De la Fuente, \emph{Mathematical Methods and Models for Economists}}{%
1 & Review of Basic Concepts (1.1--1.4, 1.6) & 1.1, 1.2 \\
2 & Metric and Normed Spaces & 3.1--3.5, 4.1 \\
3 & Vector Spaces and Linear Transformations & 2.1, 2.2, 2.4 \\
4 & Differential Calculus & 5.1, 5.2, 5.4 \\
5 & Static Models and Comparative Statics (5.1--5.3) & 4.3, 5.3 \\
6 & Convex Sets and Concave Functions (con apéndice) & 2.5, 6.1--6.4 \\
7 & Static Optimization (7.1--7.2) & 4.2, 7.1, 7.3--7.5, 8.1 \\
8 & Some Applications to Microeconomics (8.1--8.3) & 8.2 \\
9 & Dynamical Systems I: Basic Concepts and Scalar Systems & 9.1--9.3 \\
10 & Dynamical Systems II: Higher Dimensions (10.1--10.3) & 10.1, 10.2 \\
11 & Dynamical Systems III: Some Applications (11.1--11.4) & 10.4 \\
12 & An Introduction to Dynamic Optimization & 12.1, 12.2, 12.5, 13.1, 13.2 \\
13 & Some Applications of Dynamic Optimization (13.1--13.4) & 14.1--14.3 \\
}

\mapa{A.4 Kamien y Schwartz, \emph{Dynamic Optimization}}{%
I.1--I.18 & Calculus of Variations (completa) & 11.1--11.5 \\
II.1--II.18 & Optimal Control (hasta Jumps in the State Variable) & 12.1--12.6, 12.8 \\
II.21 & Dynamic Programming & 13.7 \\
}

\mapa{A.5 Lomelí y Rumbos, \emph{Métodos dinámicos en economía}}{%
2--3 & Ecuaciones diferenciales lineales; no lineales de primer orden & 9.2 \\
4--5 & Sistemas de ecuaciones diferenciales lineales; análisis cualitativo & 10.1, 10.2 \\
6--7 & Dinámica discreta; sistemas de ecuaciones en diferencias lineales & 9.3, 10.1 \\
8 & La mano invisible & 10.3 \\
9 & Ecuaciones en diferencias estocásticas & 9.4 \\
11 & Introducción al cálculo en variaciones & 11.1--11.4 \\
12--13 & Teoría de control; problemas con restricciones & 12.1, 12.3--12.5, 12.8 \\
14 & Elementos de programación dinámica & 13.1--13.3, 13.5 \\
}

\mapa{A.6 Sargent y Stachurski, \emph{Dynamic Programming}, vol.~I}{%
1 & Introduction & 13.1, 13.2 \\
2 & Operators and Fixed Points & 4.4, 13.4 \\
3 & Markov Dynamics & 13.5, 14.1 \\
4 & Optimal Stopping & 13.6 \\
5 & Markov Decision Processes (5.1--5.2) & 13.5, 14.3 \\
6 & Stochastic Discounting (6.1--6.3) & 14.4 \\
7 & Nonlinear Valuation (7.1--7.3) & 14.4, 14.5 \\
8--9 & Recursive Decision Processes; Abstract Dynamic Programming & 14.5 \\
10 & Continuous Time & 14.6 \\
}

\mapa{A.7 Meinsma y van der Schaft, \emph{A Course on Optimal Control}}{%
1 & Calculus of Variations (1.1--1.7) & 11.1--11.3 \\
2 & Minimum Principle (2.1--2.8) & 12.1--12.3 \\
3 & Dynamic Programming (3.1--3.6) & 13.1, 13.7 \\
4 & Linear Quadratic Control (4.1--4.5) & 12.7 \\
B & Differential Equations and Lyapunov Functions (B.1--B.4) & 10.3 \\
}

\mapa{A.8 Otros textos del plan}{%
--- & Bonifaz y Lama, \emph{Optimización dinámica y teoría económica}: caps.~I--IV & 11.1--11.4, 12.1--12.4, 12.8, 13.1, 13.3 \\
--- & Stachurski, \emph{Economic Dynamics}: secc.~3.1--3.2, 4.1--4.3, 10.1, 12.1 & 3.4, 10.4, 13.5, 13.8 \\
--- & Luenberger, \emph{Optimization by Vector Space Methods}: secc.~7.1--7.6, 8.1--8.6, 9.5--9.6 & 7.6, 11.6, 12.6 \\
--- & Elaydi, \emph{An Introduction to Difference Equations}: secc.~1.1--1.3, 1.5--1.7, 3.1--3.3, 4.1--4.6 & 9.3, 10.1, 10.3 \\
--- & Heijdra, \emph{Foundations of Modern Macroeconomics}: apéndice matemático, secc.~A.6 & 10.5 \\
}

%------------------------------------------------------------------------------
\section*{Anexo B. Lo que no entra en el plan}
\addcontentsline{toc}{section}{Anexo B. Lo que no entra en el plan}

{\small\filasalternas
\noindent\begin{tabularx}{\textwidth}{>{\raggedright\arraybackslash}p{4.6cm}X}
\toprule
\rowcolor{tonoOscuro}\cab{Materia} & \cab{Capítulos, secciones y programas} \\
\midrule
Teoría de juegos & \DF{8.4}; el tema h de MAT291. Está en los Bloques 6 y 7 del sílabo de microeconomía \\
Métodos numéricos & \ST{2}, \ST{6.2--6.3}, \ST{10.2}; \DP{2.1.4} y las secciones de cómputo de \DP{1.3.2} y \DP{5.1.4}; \LU{10}; Boyd y Vandenberghe (2004), caps.~9--11; las unidades h e i de MAT693; MAT237; el curso de Childers (CMU) y Judd (1998). Corresponden a otro curso \\
Estimación estructural & \DP{5.3} (ecuaciones de Bellman modificadas, \emph{Q-factors}); corresponde al sílabo de microeconometría \\
Microeconomía aplicada & \DF{8.5} (competencia imperfecta); la mayor parte de MAT656. Están en el sílabo de microeconomía \\
Teoría de la medida y dinámica estocástica avanzada & \ST{7--9}, \ST{11}, \ST{12.2}; Corbae et al.\ (2009), cap.~7 (Anexo D.1) \\
Control avanzado & \KS{II.19--II.20} (retardos y ecuaciones de estado integrales), \KS{II.22--II.23} (control estocástico y juegos diferenciales, Anexos D.2 y D.4); \MS{5}; \EL{10} \\
Ecuaciones en diferencias avanzadas & \EL{5--9} \\
Análisis funcional & \LU{2--6} (Anexo D.5) \\
Programación dinámica a nivel de posgrado en macroeconomía & Stokey, Lucas y Prescott (1989), caps.~4--5 y 9; Ljungqvist y Sargent (2018). Están en el Bloque 14 del sílabo de macroeconomía \\
\bottomrule
\end{tabularx}
}

\medskip
\noindent\destacado{Archivos de la carpeta que no se usaron por ser imágenes sin
texto}: Chiang, \emph{Elements of Dynamic Optimization}
(\texttt{elements-of-dynamic-optimization.pdf}); Cerdá, \emph{Optimización
dinámica} (\texttt{optimizacion-dinamica-2-pdf-cerda.pdf}); Dixit,
\emph{Optimization in Economic Theory}; el archivo
\texttt{Optimizacionestaticaydinamicaeneconomia.pdf}; Sargent, \emph{Dynamic
Macroeconomic Theory}; y la copia de Stokey, Lucas y Prescott de la carpeta
de macroeconomía. Cerdá, Chiang y Stokey et al.\ se citan como consulta
porque los usan MAT291, 1ECO75 y MAT616.

\medskip
\noindent\destacado{Duplicados}: Lomelí y Rumbos (en \texttt{Optimizacion} y en
\texttt{Metodos numericos}), Bonifaz y Winkelried, Chiang (\emph{Métodos
fundamentales}), MAT835 y MAT659 aparecen dos veces en la carpeta.

%------------------------------------------------------------------------------
\section*{Anexo C. Programas utilizados}
\addcontentsline{toc}{section}{Anexo C. Programas utilizados}

{\small\filasalternas
\noindent\begin{tabularx}{\textwidth}{>{\raggedright\arraybackslash}p{5.6cm}X}
\toprule
\rowcolor{tonoOscuro}\cab{Programa} & \cab{Temas del plan} \\
\midrule
MAT291, Matemáticas para economistas (pregrado, 2025-2, J.~Chávez; jefes de práctica M.~Gallardo y R.~Crousillat) & Preferencias (8.2), convexidad y optimización estática (Bloques 6 y 7), sistemas dinámicos (Bloques 9 y 10), cálculo de variaciones y control (Bloques 11 y 12), programación dinámica (Bloque 13). La teoría de juegos queda fuera \\
1ECO75, Economía Matemática Intermedia (Maestría en Economía, 2025-1, R.~Huamán y L.~Ramos) & Módulo 1, Bloque 2; módulo 2, Bloques 3, 5 y 6; módulo 3, Bloque 7; módulo 4, Bloques 10 y 12 \\
MAT693, Concavidad y Optimización (Maestría en Matemáticas Aplicadas, 2025-2, J.~E.~Martínez-Legaz y A.~Jordán) & Unidades a y b, Bloque 6; c y d, 7.2--7.5; e, 7.6; f, 8.1; g, 6.4 \\
MAT659, Temas Avanzados en Economía Matemática (2021-0, A.~Jordán) & Órdenes, 1.2; teorema del máximo, 4.2; supermodularidad y estática comparativa, 4.4 y 8.1; hipertopologías y convexidad discreta, Anexo D.3 \\
MAT835, Control Óptimo Estocástico (2025-2, J.~Farfán) & Control determinístico y ecuación HJB, 13.7; parada óptima, 13.6; tiempo continuo, 14.6; el resto, Anexo D.2 \\
MAT616, Macroeconomía Dinámica (2023-2, J.~C.~Aquino) & De referencia para el Bloque 13; su desarrollo está en el Bloque 14 del sílabo de macroeconomía \\
MAT656, Introducción a la Economía Matemática (2026-1, A.~Jordán) & De referencia para 8.2; su contenido está en el sílabo de microeconomía \\
\bottomrule
\end{tabularx}
}

\medskip
\noindent\destacado{Programas de la carpeta que no se usaron}: Investigación de
Operaciones (IOP222, IOP224 e Investigación de Operaciones 2), por ser de
ingeniería; Cálculo Numérico (MAT237) y \emph{Computational Methods for
Economics} (CMU), por ser métodos numéricos. También se revisó el
\emph{Math Camp} de Murphy (Pennsylvania): su parte de optimización coincide
con el Bloque 7 y su parte de medida y probabilidad corresponde al sílabo de
econometría.

%------------------------------------------------------------------------------
\section*{Anexo D. Material opcional}
\addcontentsline{toc}{section}{Anexo D. Material opcional}

Todo lo de este anexo es opcional: no entra en el cómputo de horas ni modifica
el cronograma.

\subsection*{D.1 Teoría de la medida y dinámica estocástica}
\ST{7--9} (integración, cadenas de Markov con densidad y probabilidad con
teoría de la medida), \ST{11} (dinámica estocástica) y \ST{12.2} (recompensas
no acotadas); Corbae et al.\ (2009), cap.~7. Momento sugerido: con el tema
13.8.

\subsection*{D.2 Control óptimo estocástico}
Sigue el programa MAT835: \KS{II.22} como introducción; después Øksendal
(2010), Fleming y Soner (2010) y Pham (2010), que no están en la carpeta.
Requiere D.1 y el tema 14.6.

\subsection*{D.3 Análisis variacional y convexidad discreta}
Sigue la segunda y la tercera parte de MAT659: convergencia de conjuntos y
estabilidad de problemas de optimización (Rockafellar y Wets, 2009; Royset,
2020) y convexidad discreta con bienes indivisibles (Murota, 1998; Danilov,
Koshevoy y Murota, 2001).

\subsection*{D.4 Juegos diferenciales y control con retardos}
\KS{II.19--II.20} y \KS{II.23}. Requiere el Bloque 12 y los juegos dinámicos
del sílabo de microeconomía.

\subsection*{D.5 Optimización en espacios vectoriales}
\LU{2--6}: espacios de Hilbert, espacios duales, Hahn--Banach y operadores
adjuntos. Es la base funcional de los temas 7.6, 11.6 y 12.6.

\subsection*{D.6 Optimización convexa}
Boyd y Vandenberghe (2004), caps.~2--5, para profundizar los Bloques 6 y 7 con
la teoría de dualidad y los teoremas de la alternativa (secc.~5.8).

\subsection*{D.7 Sistemas dinámicos en castellano}
Chávez, \emph{Sistemas dinámicos para el análisis económico}, de acceso libre
en el repositorio institucional de la PUCP: es el complemento natural de
Chávez y Gallardo para los Bloques 9 y 10. Bonifaz y Winkelried (2003)
cubre sucesiones, series, ecuaciones diferenciales y en diferencias con
aplicaciones peruanas.

%------------------------------------------------------------------------------
\section*{Anexo E. Referencias}
\addcontentsline{toc}{section}{Anexo E. Referencias}

\subsection*{E.1 Textos del plan}
{\small
\begin{itemize}[leftmargin=1.4em,itemsep=1pt,label={}]
\item Bonifaz, J. L. y R. Lama (2013). \emph{Optimización dinámica y teoría económica}, 1.ª ed.\ corregida. Apuntes de Estudio 33. Universidad del Pacífico.
\item Carter, M. (2001). \emph{Foundations of Mathematical Economics}. The MIT Press.
\item Chávez, J. y M. Gallardo (2026). \emph{Álgebra lineal y optimización para el análisis económico}. Fondo Editorial de la PUCP.
\item De la Fuente, A. (2000). \emph{Mathematical Methods and Models for Economists}. Cambridge University Press.
\item Elaydi, S. (2005). \emph{An Introduction to Difference Equations}, 3.ª ed. Springer.
\item Heijdra, B. J. (2017). \emph{Foundations of Modern Macroeconomics}, 3.ª ed., con L. S. M. Reijnders y W. E. Romp. Oxford University Press (carpeta de macroeconomía).
\item Kamien, M. I. y N. L. Schwartz (1991). \emph{Dynamic Optimization: The Calculus of Variations and Optimal Control in Economics and Management}, 2.ª ed. Elsevier.
\item Lomelí, H. y B. Rumbos. \emph{Métodos dinámicos en economía: otra búsqueda del tiempo perdido}. ITAM. Versión de la carpeta (2001); publicado por Thomson en 2003.
\item Luenberger, D. G. (1969). \emph{Optimization by Vector Space Methods}. Wiley.
\item Meinsma, G. y A. van der Schaft (2023). \emph{A Course on Optimal Control}. Springer.
\item Sargent, T. J. y J. Stachurski (2025). \emph{Dynamic Programming, Volume I: Finite States}. Versión de los autores, abril de 2025.
\item Stachurski, J. (2022). \emph{Economic Dynamics: Theory and Computation}, 2.ª ed. The MIT Press. Versión preliminar de la carpeta (2024).
\end{itemize}
}

\subsection*{E.2 Libros de consulta de la carpeta (sin horas)}
{\small
\begin{itemize}[leftmargin=1.4em,itemsep=1pt,label={}]
\item Acemoglu, D. (2009). \emph{Introduction to Modern Economic Growth}. Princeton University Press (carpeta de desarrollo y crecimiento).
\item Bonifaz, J. L. y D. Winkelried (2003). \emph{Matemáticas para la economía dinámica}, 1.ª ed.\ corregida. Universidad del Pacífico.
\item Boyd, S. y L. Vandenberghe (2004). \emph{Convex Optimization}. Cambridge University Press.
\item Corbae, D., M. B. Stinchcombe y J. Zeman (2009). \emph{An Introduction to Mathematical Analysis for Economic Theory and Econometrics}. Princeton University Press.
\item Simon, C. P. y L. Blume (1994). \emph{Mathematics for Economists}. W. W. Norton.
\end{itemize}
}

\subsection*{E.3 Referencias tomadas de los programas}
{\small
\begin{itemize}[leftmargin=1.4em,itemsep=1pt,label={}]
\item Bridges, D. y G. Mehta (1995). \emph{Representations of Preference Orderings}. Lecture Notes in Economics and Mathematical Systems. Springer.
\item Cerdá, E. (2001). \emph{Optimización dinámica}. Prentice Hall.
\item Chiang, A. C. (1992). \emph{Elements of Dynamic Optimization}. McGraw-Hill.
\item Danilov, V., G. Koshevoy y K. Murota (2001). Discrete Convexity and Equilibria in Economies with Indivisible Goods and Money. \emph{Mathematical Social Sciences}, 41, 251--273.
\item Fleming, W. H. y H. M. Soner (2010). \emph{Controlled Markov Processes and Viscosity Solutions}, 2.ª ed. Springer.
\item Léonard, D. y N. Van Long (1992). \emph{Optimal Control Theory and Static Optimization in Economics}. Cambridge University Press.
\item Murota, K. (1998). Discrete Convex Analysis. \emph{Mathematical Programming}, 83, 313--371.
\item Øksendal, B. (2010). \emph{Stochastic Differential Equations: An Introduction with Applications}, 6.ª ed. Springer.
\item Pham, H. (2010). \emph{Continuous-time Stochastic Control and Optimization with Financial Applications}. Springer.
\item Rockafellar, R. T. y R. Wets (2009). \emph{Variational Analysis}. Springer.
\item Royset, J. O. (2020). Set-Convergence and Its Applications: A Tutorial. \emph{Set-Valued and Variational Analysis}, 28, 707--732.
\item Stokey, N. L. y R. E. Lucas, con E. C. Prescott (1989). \emph{Recursive Methods in Economic Dynamics}. Harvard University Press.
\item Topkis, D. M. (1998). \emph{Supermodularity and Complementarity}. Princeton University Press.
\end{itemize}
}

\subsection*{E.4 Otras referencias}
{\small
\begin{itemize}[leftmargin=1.4em,itemsep=1pt,label={}]
\item Chávez, J. \emph{Sistemas dinámicos para el análisis económico}. Repositorio institucional de la PUCP (acceso libre).
\item Judd, K. L. (1998). \emph{Numerical Methods in Economics}. The MIT Press (citado solo en el Anexo B).
\item Ljungqvist, L. y T. J. Sargent (2018). \emph{Recursive Macroeconomic Theory}, 4.ª ed. The MIT Press (citado solo en el Anexo B).
\end{itemize}
}

\vspace{1.2cm}
\franjas[2pt]
\begin{center}\small\color{tonoPizarra}
Plan de lectura de elaboración propia. Los títulos de los capítulos
corresponden a las obras citadas y se reproducen únicamente con fines de
referencia bibliográfica. La asignación de capítulos y secciones se basa en
los índices de los libros de la carpeta y, para Chávez y Gallardo, en el
índice del ejemplar físico; los temas de los cursos proceden de los
programas citados en el Anexo C.
\end{center}

\end{document}
